% ============================================================
%  CHƯƠNG 3: MÔ HÌNH HỖN HỢP GAUSSIAN (GMM)
%  Người viết: [Họ tên thành viên phụ trách]
% ============================================================

\section{Mô hình Hỗn hợp Gaussian}
\label{sec:gmm}

% -----------------------------------------------------------
\subsection{Trực giác và động lực}
% -----------------------------------------------------------
% Gợi ý nội dung:
%   - Tại sao một Gaussian đơn không đủ mô tả dữ liệu phức tạp?
%   - Ý tưởng: dùng tổng có trọng số của nhiều Gaussian (hỗn hợp)
%   - Ví dụ trực quan: phân phối chiều cao nam/nữ, màu pixel trong ảnh

% TODO: Viết nội dung ở đây


% -----------------------------------------------------------
\subsection{Định nghĩa mô hình GMM}
% -----------------------------------------------------------
% Gợi ý nội dung:
%   - Công thức: p(x|θ) = Σ_k π_k N(x | μ_k, Σ_k)
%   - Bộ tham số θ = {π_k, μ_k, Σ_k}_{k=1}^K
%   - Ràng buộc: π_k ≥ 0, Σ_k π_k = 1, Σ_k xác định dương
%   - K là siêu tham số (hyperparameter)

% TODO: Viết nội dung ở đây


% -----------------------------------------------------------
\subsection{Mô hình sinh (Generative model) với biến ẩn}
% -----------------------------------------------------------
% Gợi ý nội dung:
%   - Biến ẩn one-hot z = (z_1,…,z_K): z_k = 1 nghĩa là chọn thành phần k
%   - Phân phối tiên nghiệm: p(z_k=1) = π_k  →  Categorical(π)
%   - Phân phối có điều kiện: p(x | z_k=1, θ) = N(x | μ_k, Σ_k)
%   - Phân phối kết hợp: p(x,z|θ) = Π_k [π_k N(x|μ_k,Σ_k)]^{z_k}
%   - Sơ đồ plate notation (nếu có, vẽ bằng TikZ)

% TODO: Viết nội dung ở đây

% Ví dụ sơ đồ plate notation (bỏ comment khi vẽ xong):
% \begin{figure}[H]
%   \centering
%   \input{figures/plate_notation.tex}   % hoặc \includegraphics
%   \caption{Sơ đồ plate notation của GMM}
%   \label{fig:plate}
% \end{figure}


% -----------------------------------------------------------
\subsection{Hàm log-likelihood của dữ liệu quan sát}
% -----------------------------------------------------------
% Gợi ý nội dung:
%   - ℓ(θ) = Σ_n log [Σ_k π_k N(x_n | μ_k, Σ_k)]
%   - Vấn đề: log của tổng → không có nghiệm dạng đóng
%   - Vì sao không tối ưu trực tiếp? → cần thuật toán EM

% TODO: Viết nội dung ở đây


% -----------------------------------------------------------
\subsection{Tính chất và hạn chế của GMM}
% -----------------------------------------------------------
% Gợi ý nội dung:
%   - Định lý xấp xỉ phổ dụng: GMM xấp xỉ được mọi phân phối liên tục
%   - Vấn đề singularity (Σ_k suy biến)
%   - Sự nhạy cảm với khởi tạo và số thành phần K
%   - Lựa chọn K: BIC, AIC, cross-validation

% TODO: Viết nội dung ở đây
